\documentclass[11pt]{article}
\usepackage[margin=1in]{geometry}
\usepackage{amsmath,amssymb,amsthm}
\usepackage{graphicx}
\usepackage[hidelinks]{hyperref}

\newtheorem{theorem}{Theorem}
\newtheorem{lemma}{Lemma}
\newtheorem{proposition}{Proposition}
\newtheorem{remark}{Remark}

\newcommand{\C}{\mathbb C}
\newcommand{\N}{\mathbb N}
\newcommand{\Sinf}{S_\infty}
\newcommand{\norm}[1]{\left\lVert #1\right\rVert}

\title{Ray's \(B_{p,q}\) Isometric Embedding Conjecture}
\author{FA/Banach discovery run \texttt{fa\_banach\_001}}
\date{June 13, 2026}

\begin{document}
\maketitle

\section{Source Question}

The source paper is Samya Kumar Ray, ``On Isometric Embedding
\(\ell_p^m\to S_\infty^n\) and Unique operator space structure'',
arXiv:1911.00241.  For \(1\leq p,q<\infty\), Ray defines
\[
  B_{p,q}=\{(z_1,z_2)\in\C^2: |z_1|^p+|z_2|^q<1\}
\]
and writes \(\|\cdot\|_{B_{p,q}}\) for the norm whose open unit ball is
\(B_{p,q}\).  In Remark 4.2, after proving non-embedding for
\(2<p,q<\infty\) and giving a triangular positive example, Ray conjectures
that
\[
(\C^2,\|\cdot\|_{B_{p,q}})\hookrightarrow \Sinf
\quad\text{isometrically}
\]
if and only if
\[
  (p,q)=(2,2),\qquad (p,q)=(1,2),\qquad\text{or}\qquad (p,q)=(2,1).
\]

\begin{center}
\includegraphics[width=\textwidth]{figures/open_problem_crop.png}
\end{center}

\section{The Result}

\begin{theorem}
Let \(1\leq p,q<\infty\).  The two-dimensional complex Banach space
\((\C^2,\|\cdot\|_{B_{p,q}})\) embeds linearly isometrically into
\(\Sinf\) if and only if
\[
  (p,q)\in\{(2,2),(1,2),(2,1)\}.
\]
\end{theorem}

\section{Positive Cases}

For \((p,q)=(2,2)\), the map
\[
  (z_1,z_2)\longmapsto
  \begin{pmatrix} z_1 & z_2\\ 0&0\end{pmatrix}
\]
has operator norm \((|z_1|^2+|z_2|^2)^{1/2}\).

For the mixed cases, use the elementary identity
\[
  \left\|
  \begin{pmatrix} b&a\\0&b\end{pmatrix}
  \right\|\leq 1
  \quad\Longleftrightarrow\quad
  |a|+|b|^2\leq 1.
\]
Indeed, by rotation invariance assume \(a,b\geq0\).  The largest eigenvalue
of
\[
\begin{pmatrix} b&a\\0&b\end{pmatrix}^*
\begin{pmatrix} b&a\\0&b\end{pmatrix}
\]
is \(1\) exactly when \(a+b^2=1\), as follows by checking the determinant of
\(A^*A-I\), and monotonicity gives the inequality case.  Hence
\[
  (z_1,z_2)\longmapsto z_1I_2+z_2e_{12}
\]
is an isometry for \(B_{2,1}\), while
\[
  (z_1,z_2)\longmapsto z_2I_2+z_1e_{12}
\]
is an isometry for \(B_{1,2}\).  These matrices are finite-rank compact
operators, so the range lies in \(\Sinf\).

\section{Two Obstructions}

We use two negative ingredients.  The first is Ray's Theorem 1.1:
complex \(\ell_1^2\) does not embed linearly isometrically into \(\Sinf\).
This excludes \(B_{1,1}\).

The second is the already-packaged partial result
\[
  \text{if } \max\{p,q\}>2,\quad
  (\C^2,\|\cdot\|_{B_{p,q}})\not\hookrightarrow \Sinf.
\]
For completeness, recall the short reason.  Ray's Remark 4.2 gives a
criterion reducing non-embedding to the fact that for every \(0<c\leq1\),
the map \((z_1,z_2)\mapsto(z_1,cz_2)\) is not a contraction into
\(\ell_2^2\).  If \(q>2\), choose \(y>0\) small and
\[
  x=(1-y^q)^{1/p}.
\]
Then \((x,y)\) is on the \(B_{p,q}\)-unit sphere and
\[
  x^2+c^2y^2=(1-y^q)^{2/p}+c^2y^2>1
\]
for all sufficiently small \(y\), since \(y^2\) dominates \(y^q\).  Thus the
criterion applies.  The case \(p>2\) follows by swapping coordinates.

It remains only to exclude the fractional region where one exponent lies
strictly between \(1\) and \(2\).

\section{Axis-Curve Lemma}

\begin{lemma}
Let \(T,S\in \Sinf\) and \(\|T\|=1\).  For \(t\geq0\) near \(0\), the
function
\[
  g(t)=\|T+tS\|^2
\]
has the following property: if \(g(t)>1\) for all sufficiently small
positive \(t\), then
\[
  g(t)-1 = c t^m + o(t^m)
\]
for some integer \(m\geq1\) and some \(c>0\).
\end{lemma}

\begin{proof}
Set
\[
  A(t)=(T+tS)^*(T+tS)
      =T^*T+t(T^*S+S^*T)+t^2S^*S.
\]
This is a self-adjoint compact analytic family.  Since \(T\) is compact and
\(\|T\|=1\), the eigenvalue \(1\) of \(A(0)=T^*T\) has finite multiplicity
and is isolated from the rest of the spectrum.  Let \(\Gamma\) be a small
circle enclosing \(1\) and no other point of \(\sigma(A(0))\).  For small
\(|t|\), the Riesz projection
\[
  P(t)=\frac{1}{2\pi i}\int_\Gamma (z-A(t))^{-1}\,dz
\]
is finite-rank and analytic in \(t\), and the spectral values of \(A(t)\)
inside \(\Gamma\) are precisely the eigenvalues of a finite-dimensional
self-adjoint analytic family.  By Rellich's theorem, these eigenvalues can
be written as finitely many real analytic functions
\[
  \lambda_1(t),\ldots,\lambda_n(t),\qquad \lambda_j(0)=1.
\]
The rest of the spectrum of \(A(t)\) stays below \(1-\delta\) for some
\(\delta>0\).  Hence, for small \(t\geq0\),
\[
  g(t)=\max_{1\leq j\leq n}\lambda_j(t).
\]
Each \(\lambda_j(t)-1\) has an ordinary Taylor expansion with integer
powers.  Since \(g(t)>1\) for small positive \(t\), at least one branch has
a first nonzero Taylor coefficient that is positive on the positive side.
Taking the smallest such leading order among the dominating branches gives
the claimed integer \(m\) and coefficient \(c>0\).
\end{proof}

\begin{lemma}
If \(1<q<2\), then
\((\C^2,\|\cdot\|_{B_{p,q}})\) does not embed linearly isometrically into
\(\Sinf\).
\end{lemma}

\begin{proof}
Suppose that \(\Phi(z_1,z_2)=z_1T+z_2S\) is a linear isometry into
\(\Sinf\).  Then \(\|T\|=1\) and, for \(t\geq0\),
\[
  \|T+tS\|=\|(1,t)\|_{B_{p,q}}.
\]
Let \(\lambda(t)=\|(1,t)\|_{B_{p,q}}\).  By definition,
\[
  \lambda(t)^{-p}+t^q\lambda(t)^{-q}=1.
\]
Writing \(\lambda(t)=1+u(t)\) and using
\((1+u)^{-p}=1-pu+O(u^2)\), we get
\[
  u(t)=\frac{1}{p}t^q+o(t^q).
\]
Thus
\[
  \|T+tS\|^2-1=\lambda(t)^2-1
  =\frac{2}{p}t^q+o(t^q).
\]
This contradicts the axis-curve lemma, because \(q\in(1,2)\) is not an
integer.  Therefore no isometric embedding exists.
\end{proof}

\begin{lemma}
If \(1<p<2\), then
\((\C^2,\|\cdot\|_{B_{p,q}})\) does not embed linearly isometrically into
\(\Sinf\).
\end{lemma}

\begin{proof}
Apply the previous lemma after swapping the two coordinates, or repeat the
same argument along the second coordinate axis.  Explicitly,
\(\|(t,1)\|_{B_{p,q}}=1+(1/q)t^p+o(t^p)\), which is incompatible with the
integer-power conclusion of the axis-curve lemma when \(1<p<2\).
\end{proof}

\section{Conclusion}

If \(\max\{p,q\}>2\), non-embedding follows from Ray's descent criterion as
recalled above.  If \(1<p<2\) or \(1<q<2\), non-embedding follows from the
axis-curve lemma.  The remaining possibilities are
\[
  p,q\in\{1,2\}.
\]
Among these four pairs, \((1,1)\) is excluded by Ray's complex
\(\ell_1^2\) theorem, and \((2,2),(1,2),(2,1)\) have the explicit finite-rank
embeddings above.  This proves the theorem.

\section{Verification and Novelty Check}

The script \texttt{code/check\_axis\_and\_embeddings.py} numerically checks
the axis asymptotics and the triangular positive embedding.  These checks are
not used as proof.

Bounded novelty check on June 13, 2026: exact web searches included
``\(B_{p,q}\)'', ``\(S_\infty\)'', ``embeds isometrically'', ``Ray'',
``conjecture'', and ``unique operator space structure'' in the relevant
combinations.  These searches surfaced Ray's original arXiv page but no later
resolution of the conjecture.  The local run registry had only the previous
\(\max\{p,q\}>2\) partial packet.

Human review should focus on the use of Rellich's theorem in the axis-curve
lemma and on the asymptotic comparison
\(\|(1,t)\|_{B_{p,q}}=1+(1/p)t^q+o(t^q)\).

\end{document}
